\section{时间—地理—产品—客户—物流的双证据驱动识别}

问题2要区分“组间销售分布有差异”与“该因素能解释未来订单”这两个命题。治理后的样本含9799条明细，% src:求解/问题2/结果/运行摘要.json:清洗后明细数
但只对应4921个订单，% src:求解/问题2/结果/运行摘要.json:订单数
同一订单中的商品行并不独立。因为本题同时存在强右尾和订单内相关，所以我们保留商品明细以免抹去产品结构，又把检验、重采样和验证损失统一提升到订单簇层面。五维字段进入两条相互独立的证据链：秩检验回答分布差异，时间外置换回答预测增量，跨年方向和参数翻转审查负责拦截偶然信号。

\subsection{订单簇稳健检验与时间外预测证据}

销售额的稳健尺度定义为
\begin{equation}
    Y_i=\log(1+S_i),
    \label{eq:q2_log_sales}
\end{equation}
其中，$S_i$为第$i$条商品明细的销售额，$Y_i$为其对数销售额。该变换不改变秩次，却压缩大额明细对预测损失的支配程度。对任一含$k$个水平的因素$F$，以平均秩偏离构造Kruskal--Wallis统计量\cite{kruskal1952ranks}，并用epsilon平方描述差异规模：
\begin{equation}
    H_F=\frac{12}{n(n+1)}\sum_{g=1}^{k}n_g
    \left(\bar R_g-\frac{n+1}{2}\right)^2,
    \qquad
    \varepsilon_F^2=\max\left\{0,\frac{H_F-k+1}{n-k}\right\}.
    \label{eq:q2_rank_effect}
\end{equation}
其中，$n_g$和$\bar R_g$分别表示第$g$组的明细数与平均秩，$n$为总明细数。普通渐近检验会把每条商品行视为独立观测，本题则只在明细数相同的订单之间交换完整标签块，并按订单整簇bootstrap；多个因素的显著性再用BH方法控制假发现率\cite{benjamini1995fdr}。该设计来自一单多行的实际结构，而不是对普通检验作形式替换。

客户和产品编号具有高基数，直接独热编码会形成稀疏特征，全样本目标编码还会把当前销售额写回自身。对实体$e$在日期$t_i$出现的明细，本文仅使用严格早于该日的历史记录，定义平滑历史价值
\begin{equation}
    \widetilde{\mu}_{e,i}
    =\frac{\sum_{j:e_j=e,\,t_j<t_i}Y_j
    +\lambda\mu_{<t_i}}
    {\left|\{j:e_j=e,\,t_j<t_i\}\right|+\lambda}.
    \label{eq:q2_history_encoding}
\end{equation}
其中，$\mu_{<t_i}$为$t_i$之前的全局对数销售均值，$\lambda$为向全局历史水平收缩的先验强度。严格使用$t_j<t_i$，使同日其他商品和未来年度信息均不能进入当前特征；产品历史价值层的分位边界也只由训练期估计。

第二条证据采用绝对误差损失的低自由度梯度提升树\cite{friedman2001gbm}。训练期含6541条明细和3260个订单，% src:求解/问题2/结果/运行摘要.json:时间外模型性能.训练明细数、训练订单数
时间外验证期含3258条明细和1661个订单，% src:求解/问题2/结果/运行摘要.json:时间外模型性能.测试明细数、测试订单数
两个时期的订单编号没有交叉。2025年的角色是“被81组翻转诊断重复使用的时间外验证集”，% src:求解/问题2/结果/关键参数灵敏度.json:重要性方向判定、组合总数
只用于检查因素身份和重要性方向是否翻转，不据其误差另选主设定。严格历史特征按订单到达逐日更新；若经营场景要求在年初一次性预测全年，需另行固定历史截止日验证，不能直接沿用该逐日可得口径。设验证订单$c$含$m_c$条明细，因素组$G$被整簇置换后的预测为$\widehat Y_i^{\pi_G}$，则其时间外重要性为
\begin{equation}
    \Delta_G=\frac{1}{C}\sum_{c=1}^{C}\frac{1}{m_c}
    \sum_{i\in\mathcal O_c}
    \left(\left|Y_i-\widehat Y_i^{\pi_G}\right|
    -\left|Y_i-\widehat Y_i\right|\right).
    \label{eq:q2_permutation_importance}
\end{equation}
其中，$\mathcal O_c$为订单$c$的商品行集合，$C$为验证订单数。$\Delta_G>0$表示破坏该因素后订单等权误差上升；先在订单内平均再对订单平均，可防止多商品订单自动获得更高权重。实际履约天数发生在下单之后，因此不进入预测模型，物流维只保留下单时已知的发货模式。

图\ref{fig:q2_path}连接了五维字段、两条证据、跨年检验与参数审查。因素需同时满足FDR统计证据、稳定正的时间外重要性和跨年方向要求，基准设置入选后还要在全部参数组合中不翻转，才被认定为稳定核心。

\begin{figure}[H]
  \centering
  \includegraphics[width=0.8\textwidth]{../求解/问题2/图片/五维因素作用路径图.png}
  \caption{五维因素的双证据筛选与经营使用路径}
  \label{fig:q2_path}
\end{figure}

\subsection{产品因素的双证据强度}

图\ref{fig:q2_effect_heatmap}把效应量、FDR显著性与最终层级放在同一尺度上。产品历史价值层的效应量最高，产品子类次之，产品大类仍保留清晰的组间秩差；区域虽然达到统计显著，但效应量接近于零，显著性并不等同于经营影响强。

\begin{figure}[H]
  \centering
  \includegraphics[width=0.8\textwidth]{../求解/问题2/图片/因素效应量与显著性热图.png}
  \caption{五维因素的效应量、显著性与结论层级}
  \label{fig:q2_effect_heatmap}
\end{figure}

\begingroup
\small
\setlength{\tabcolsep}{5pt}
\begin{longtable}{>{\centering\arraybackslash}p{0.15\textwidth}>{\centering\arraybackslash}p{0.16\textwidth}>{\centering\arraybackslash}p{0.17\textwidth}>{\centering\arraybackslash}p{0.18\textwidth}>{\centering\arraybackslash}p{0.18\textwidth}}
\caption{核心产品因素与区域信号的证据对照}\label{tab:q2_evidence}\\
\toprule
因素 & $\varepsilon^2$ & $\Delta\mathrm{MAE}$ & 跨年同向率 & 网格入选 \\
\midrule
\endfirsthead
\multicolumn{5}{c}{续表~\thetable}\\
\toprule
因素 & $\varepsilon^2$ & $\Delta\mathrm{MAE}$ & 跨年同向率 & 网格入选 \\
\midrule
\endhead
\midrule
\multicolumn{5}{r}{续下页}\\
\endfoot
\bottomrule
\endlastfoot
产品历史价值层 &
$0.5630$% src:求解/问题2/结果/稳健组间检验.json:结果[因素=产品历史价值层].epsilon平方效应量
& $0.9988$% src:求解/问题2/结果/时间外分组置换重要性.json:重要性结果[因素或维度=产品历史].置换后订单等权log1p_MAE增量
& $1.00$% src:求解/问题2/结果/跨年稳定性.json:结果[因素=产品历史价值层].四年方向一致率
& $81/81$% src:求解/问题2/结果/关键参数灵敏度.json:各因素入选次数.产品历史价值层、组合总数
\\
产品子类 &
$0.4432$% src:求解/问题2/结果/稳健组间检验.json:结果[因素=产品子类].epsilon平方效应量
& $0.0301$% src:求解/问题2/结果/时间外分组置换重要性.json:重要性结果[因素或维度=产品子类].置换后订单等权log1p_MAE增量
& $1.00$% src:求解/问题2/结果/跨年稳定性.json:结果[因素=产品子类].四年方向一致率
& $81/81$% src:求解/问题2/结果/关键参数灵敏度.json:各因素入选次数.产品子类、组合总数
\\
产品大类 &
$0.2252$% src:求解/问题2/结果/稳健组间检验.json:结果[因素=产品大类].epsilon平方效应量
& $0.0258$% src:求解/问题2/结果/时间外分组置换重要性.json:重要性结果[因素或维度=产品大类].置换后订单等权log1p_MAE增量
& $1.00$% src:求解/问题2/结果/跨年稳定性.json:结果[因素=产品大类].四年方向一致率
& $81/81$% src:求解/问题2/结果/关键参数灵敏度.json:各因素入选次数.产品大类、组合总数
\\
区域 &
$0.0023$% src:求解/问题2/结果/稳健组间检验.json:结果[因素=区域].epsilon平方效应量
& $0.0005$% src:求解/问题2/结果/时间外分组置换重要性.json:重要性结果[因素或维度=区域].置换后订单等权log1p_MAE增量
& $0.75$% src:求解/问题2/结果/跨年稳定性.json:结果[因素=区域].四年方向一致率
& $54/81$% src:求解/问题2/结果/关键参数灵敏度.json:各因素入选次数.区域、组合总数
\\
\end{longtable}
\endgroup

表\ref{tab:q2_evidence}显示，产品维内部也有明显层次。产品历史价值层的时间外增量远高于类别标签，说明同一产品过去形成的销售层级比静态类别更能区分明细销售额；子类效应又高于大类，表明具体产品结构包含粗粒度分类未表达的信息。整组置换产品维后，订单等权对数MAE增加1.1092，% src:求解/问题2/结果/时间外分组置换重要性.json:重要性结果[因素或维度=产品维].置换后订单等权log1p_MAE增量
其订单簇bootstrap区间为$[1.0837,1.1326]$，% src:求解/问题2/结果/时间外分组置换重要性.json:重要性结果[因素或维度=产品维].订单簇Bootstrap_95%区间
产品信息因而构成五维解释中的主轴。

同口径朴素基线只用训练期普通ANOVA筛选和分类均值预测，在相同验证订单上的MAE为1.1059；% src:求解/问题2/结果/基线对比.json:计划内同口径朴素基线.朴素方法2025订单等权log1p_MAE
双证据模型降至0.5904，% src:求解/问题2/结果/基线对比.json:计划内同口径朴素基线.主方法2025订单等权log1p_MAE
相对改善46.61\%。% src:求解/问题2/结果/基线对比.json:计划内同口径朴素基线.主方法相对朴素方法MAE改善率
图\ref{fig:q2_mosaic}进一步显示三项产品因素在统计、时间外、跨年和参数审查中保持同向，而区域只在前三类证据下看似完整，参数稳健性并未填满。

\begin{figure}[H]
  \centering
  \includegraphics[width=0.8\textwidth]{../求解/问题2/图片/驱动证据一致性马赛克图.png}
  \caption{候选因素在多类证据中的一致性}
  \label{fig:q2_mosaic}
\end{figure}

\subsection{阈值翻转、经营短板与解释边界}

我们最初采用四年方向一致率不低于$0.75$的判据，% src:求解/问题2/结果/跨年稳定性.json:结果[因素=区域].方向一致率阈值
区域与三项产品因素一同进入基准清单；但把阈值提高到$1.00$后，% src:求解/问题2/结果/关键参数灵敏度.json:跨年判据网格
区域因同向率只有$0.75$而退出，% src:求解/问题2/结果/跨年稳定性.json:结果[因素=区域].四年方向一致率
核心身份发生翻转。因此返工没有选择能保留原结论的单一参数，而是把历史平滑、树复杂度和跨年判据组成81组审查：% src:求解/问题2/结果/关键参数灵敏度.json:组合总数
三项产品因素均为$81/81$入选，% src:求解/问题2/结果/关键参数灵敏度.json:各因素入选次数、组合总数
区域仅为$54/81$，% src:求解/问题2/结果/关键参数灵敏度.json:各因素入选次数.区域、组合总数
且网格中的重要性方向会改变。区域据此降为阈值敏感的次级地理信号，只能用于分层观察，不能与产品核心并列。
% process:交接/实验记录.json:问题2/正确性复核的稳定阈值灵敏度

经营矩阵把稳定因素的销售份额与同比变化交叉。图\ref{fig:q2_quadrant}中，Technology在验证年度占37.31\%且同比增长21.36\%，% src:求解/问题2/结果/增长驱动与短板矩阵.json:结果[因素=产品大类,水平=Technology].2025销售份额、2025较2024增长率
可作为产品布局中的增长承接对象；Machines份额为6.03\%，却同比下降22.11\%，% src:求解/问题2/结果/增长驱动与短板矩阵.json:结果[因素=产品子类,水平=Machines].2025销售份额、2025较2024增长率
属于规模不低但正在收缩的产品短板。West虽表现为34.36\%的份额和35.98\%的增长，% src:求解/问题2/结果/增长驱动与短板矩阵.json:结果[因素=区域,水平=West].2025销售份额、2025较2024增长率
仍只能作为容量配置线索，因为区域因素没有通过全网格核心条件。

\begin{figure}[H]
  \centering
  \includegraphics[width=0.8\textwidth]{../求解/问题2/图片/子品类—区域增长贡献象限图.png}
  \caption{子品类与区域的销售贡献—增长象限}
  \label{fig:q2_quadrant}
\end{figure}

时间、州、客户和物流标签未进入稳定核心，并不等于这些字段没有管理价值，而是现有数据不足以支持稳定的增销解释。模型在原尺度上的WAPE为77.69\%，% src:求解/问题2/结果/运行摘要.json:时间外模型性能.测试原尺度WAPE
它适合识别相对驱动和结构短板，不承担单笔金额精确预测。数据还缺少数量、折扣、价格、利润、促销和库存，故本节的“驱动”仅指稳健关联与预测信息；后续经营策略应把产品核心用于资源优先级，把区域信号用于小范围分层验证，不能把标签相关性改写成因果收益。
